\documentclass[12pt,a4paper]{article}
\usepackage[UTF8]{ctex}
\usepackage{amsmath,amssymb}
\usepackage{geometry}
\geometry{left=2.5cm,right=2.5cm,top=2.5cm,bottom=2.5cm}
\usepackage{enumitem}
\usepackage{titlesec}
\usepackage[hidelinks]{hyperref}
\usepackage{booktabs}

\title{\bfseries 文献综述：基于神经算子的托卡马克等离子体剖面拟合方法}
\author{}
\date{\today}

\begin{document}
\maketitle

\begin{abstract}
\noindent 本综述围绕「用神经网络替代传统参数化拟合，从离散诊断测量点重构连续等离子体剖面（$T_e$、$n_e$、$T_i$）」这一核心问题，梳理了三类相互关联的研究脉络：(1) 以 mtanh 为代表的传统剖面拟合范式及其局限；(2) 以 DeepONet 为核心的神经算子理论，以及 PointNet、Transformer、LSTM 等网络架构组件；(3) 深度学习在磁约束聚变诊断与剖面重建中的具体应用（RTCAKENN、NN-EFIT、EAST 快速平衡重建等）。在此基础上指出本文项目的定位：将 DeepONet 的分支--主干（branch--trunk）架构与物理约束损失相结合，从散射诊断点学习到剖面的算子映射，并嵌入 ONETWO 输运计算流程。
\end{abstract}

\section{引言}

磁约束聚变等离子体中，电子温度 $T_e$、电子密度 $n_e$ 与离子温度 $T_i$ 等径向剖面是输运分析、平衡重建与电流驱动建模的基础输入。这些剖面并非直接测量得到，而是由汤姆逊散射（TS）、反射计（Reflectometry）、X 射线晶体谱仪（XCS）等诊断系统在若干离散径向位置采样后，通过拟合与插值重建为连续曲线。剖面重建的精度与鲁棒性直接决定下游物理计算（如 ONETWO 输运求解、LHCD 电流驱动）的可靠性。

长期以来，剖面拟合依赖参数化模型与平滑插值。随着深度学习与神经算子理论的发展，一类「从离散测量到连续函数」的新范式正在兴起。本文从三条脉络梳理相关研究，并定位本文项目在其中的位置。

\section{传统剖面拟合范式及其局限}

\subsection{mtanh 台基参数化拟合}

H 模等离子体的标志性特征是在边缘形成陡峭的输运垒（edge transport barrier，ETB），对应的温度与密度剖面呈现「芯部平坦 + 台基区陡峭下降」的形态。Groebner 等（2001）系统总结了 DIII-D 装置 H 模边缘物理的定量研究，确立了用\textbf{修正双曲正切函数（mtanh）}拟合台基剖面（pedestal profile）的标准方法。mtanh 函数以芯部高度、台基顶、台基宽度、斜坡位置等少数物理参数刻画台基形态，物理可解释性强，至今仍是主流剖面分析工具。

本项目中 \texttt{fitting\_mtanh.py} 所实现的拟合链正是这一范式的典型代表：$T_e$/$n_e$ 采用「芯部样条 + mtanh 台基 + 平滑过渡」，$T_i$ 采用「mtanh + $T_e$ 台基边界约束」，并配合 IQR 迭代剔除异常点、PCHIP 保形单调插值等后处理。

\subsection{参数化方法的固有局限}

尽管 mtanh 范式物理清晰，但存在若干结构性局限：

\begin{enumerate}[label=\arabic*.]
    \item \textbf{对异常点与噪声敏感}：拟合前需人工设置阈值、迭代剔除离群点，不同炮号/时间点需反复调参，主观性强。
    \item \textbf{跨模式不鲁棒}：H 模与 L 模剖面形态差异显著，需要根据 $H_{98}$ 阈值切换不同拟合分支，模式边界附近的分类错误会直接导致拟合失败。
    \item \textbf{非物理振荡与负值}：样条/PCHIP 在数据稀疏或边缘外推时可能产生非物理振荡，甚至负值，进而导致 LHW 功率沉积积分中的 $\sqrt{\text{负值}}\to\text{NaN}$ 问题（本项目代码中对此有专门的安全裁剪防护）。
    \item \textbf{参数化表达能力受限}：固定函数形式的 mtanh 无法刻画偏离理想形态的剖面（如芯部非单调、双台基等）。
\end{enumerate}

这些局限构成了用数据驱动方法替代/辅助参数化拟合的直接动机。

\section{神经算子与 DeepONet 框架}

\subsection{从函数逼近到算子逼近}

深度神经网络是连续函数的通用逼近器，但其经典形式学习的是「向量$\to$向量」或「函数$\to$标量」的映射。许多科学与工程问题——包括从离散诊断点重构连续剖面——本质上要求学习「函数$\to$函数」的\textbf{算子映射}。Lu 等（2021）提出的\textbf{DeepONet（Deep Operator Network）}是基于算子普适逼近定理的实用实现，其核心是\textbf{双网络结构}：

\begin{itemize}
    \item \textbf{分支网络（branch net）}：在固定的传感器位置 $\{x_1,\dots,x_m\}$ 处编码输入函数，输出系数向量；
    \item \textbf{主干网络（trunk net）}：编码输出函数的位置坐标 $y$，输出一组基函数；
    \item \textbf{输出}：系数向量与基函数的点积，即 $G(u)(y) \approx \sum_k b_k(u)\cdot t_k(y)$。
\end{itemize}

这一「分支--主干点积」结构正是本项目 ProfileNet/LSTM/Transformer 三类架构共用的\textbf{SetEncoder + CoordDecoder}模式的直接理论来源：分支网络对应编码器（把散射点编码为固定维潜向量 $z$），主干网络对应 CoordDecoder（把目标 $\rho$ 坐标映射为基函数），输出 $\mathrm{Softplus}(z\cdot b(\rho))$。

DeepONet 的关键优势在于：输入函数在任意位置的传感器上采样、输出可在任意位置求值，天然适配「诊断点数量与位置可变」的剖面重构场景。

\subsection{物理信息约束的引入}

Wang 等（2021）提出\textbf{physics-informed DeepONet（PI-DeepONet）}，在 DeepONet 的数据拟合损失之外，利用自动微分施加偏微分方程的残差约束作为软惩罚，使预测结果满足底层物理规律。这一思想构成了本项目中\textbf{物理约束损失函数（PhysicsConstrainedLoss）}的直接依据：
\[
L = L_{\mathrm{MSE}} + \lambda_1 L_{\mathrm{mono}} + \lambda_2 L_{\mathrm{bdy}} + \lambda_3 L_{\mathrm{smooth}}
\]
其中 $L_{\mathrm{mono}}$ 惩罚非单调段（对应剖面径向单调递减的物理预期）、$L_{\mathrm{bdy}}$ 惩罚磁轴处非零导数（对应磁轴对称性）、$L_{\mathrm{smooth}}$ 惩罚过大二阶导数（对应剖面光滑性）。这些约束并非来自 PDE 残差，而是来自剖面先验知识，本质上是 PI 思想在「无显式 PDE、但有强物理先验」问题上的自然推广。

\subsection{神经算子概述与聚变应用}

Karniadakis 等（2021）对物理信息机器学习进行了系统综述，奠定了 PINN/神经算子的统一框架。在聚变领域，Gopakumar 等（2024）将\textbf{Fourier Neural Operator（FNO）}应用于 MAST 装置与 JOREK 模拟，展示了神经算子作为等离子体演化代理模型的潜力——相较于传统求解器获得约 6 个数量级的加速，且在归一化域内达到 $\approx 10^{-5}$ 的 MSE。该工作证明了神经算子在磁约束聚变中的可行性，但关注的是时空演化的代理建模，而非本项目聚焦的稳态径向剖面重构。

\section{网络架构组件}

本项目的四类架构（ProfileNet、LSTM、CNN、Transformer）各自对应一类基础网络原语：

\begin{itemize}
    \item \textbf{PointNet（Qi 等，2017）}：首个直接消费点云的网络，通过对称函数（max pooling）保证对输入点集排列的\textbf{置换不变性}。这是本项目\textbf{ProfileNet 的 SetEncoder}的理论基础——诊断点的测量顺序无关紧要，恰好可用 PointNet 式的对称池化编码为潜向量。
    \item \textbf{Transformer（Vaswani 等，2017）}：以多头自注意力替代循环与卷积，能建模序列中任意位置间的全局依赖。本项目\textbf{Transformer 基线}采用 2 层、4 头的 TransformerEncoder 捕获散射点间的全局相关性，与 LSTM 的递归建模形成对照。
    \item \textbf{LSTM（Hochreiter \& Schmidhuber，1997）}：长短期记忆网络通过门控机制建模序列长程依赖。本项目\textbf{BiLSTM 基线}以 $\rho$ 排序后的输入编码径向序列结构。
    \item \textbf{CNN/ResNet}：本项目\textbf{CNN-1D 基线}为纯残差卷积网络，作为不采用 DeepONet 分支--主干结构的\textbf{独立对照组}，用于验证「分支--主干点积」这一结构本身是否为性能增益的关键。
\end{itemize}

\section{深度学习在聚变诊断与剖面重建中的应用}

\subsection{实时动力学剖面重建}

Shousha 等（2024）提出的\textbf{RTCAKENN}是与本项目最直接相关的工作。该模型用机器学习近似 CAKE（一致自动动力学平衡重建）的输出，实时重构压力、反 $q$、环向电流密度、电子温度/密度、碳杂质温度与旋转等 7 条剖面，并已在 DIII-D 等离子体控制系统（PCS）中部署。其架构采用「MLP 编码标量 + 1D 卷积编码剖面输入 $\to$ MLP 潜特征提取 $\to$ 上采样 + 卷积解码」的编码器--解码器结构，且通过 dropout 训练获得在缺失 TS 或 CXRS 数据时的鲁棒性。

本项目与 RTCAKENN 的目标（从诊断测量重构剖面）高度一致，但存在三点差异：(1) 本项目直接学习「散射点$\to$剖面」的算子映射（DeepONet 式），而 RTCAKENN 采用固定网格的卷积编码；(2) 本项目显式引入物理约束损失，而非纯数据驱动；(3) 本项目针对 EAST 装置并嵌入 ONETWO/LHCD 输运流程，RTCAKENN 面向 DIII-D 实时控制。

\subsection{平衡重建的神经网络代理}

Sun 等（2024）训练多层感知机（MLP）作为 DIII-D 动力学 EFIT 的代理模型，系统研究了不同诊断输入（仅磁测量 vs.\ 加入压强/电流剖面等）对重构精度的影响，发现仅磁测量的模型已能输出较准确的动力学剖面。Lu 等（2023，arXiv:2305.12098）在 EAST 上发展了基于深度学习的快速平衡重建方法，从磁测量预测 EFIT 极向磁通，并自动优化网络超参数，达到离线 EFIT 的空间分辨率同时满足实时控制的时间约束。

这些工作表明，深度学习替代传统数值/参数化求解器已成为聚变诊断与平衡分析的重要趋势，但其重点在平衡（磁通/压强/电流）而非本项目聚焦的 $T_e$/$n_e$/$T_i$ 剖面拟合与 LHCD 输运。

\subsection{EAST 诊断基础}

Zang 等（2011）报道了 EAST 装置升级的多脉冲激光、多点汤姆逊散射（TS）诊断系统，是本项目 $T_e$ 剖面训练数据来源（TS\_EAST）的物理基础。本项目的数据集覆盖 161 炮、约 3200 个样本，$T_e$ 来自 TS，$n_e$ 来自反射计，$T_i$ 来自 XCS，正是建立在这些 EAST 诊断系统之上的。

\section{输运建模与 LHCD 背景}

剖面拟合的最终目的是为输运与电流驱动计算提供输入。本项目将拟合剖面接入 ONETWO 输运代码，并计算低杂波电流驱动（LHCD）。相关的集成建模框架（如 OMFIT，Meneghini 等，2015）与 LHCD 物理（EAST 低杂波电流驱动实验与参数衰变不稳定性研究）构成了本项目下游应用的方法学背景。与上述神经算子/实时重建文献相比，本项目的一个特色在于：不仅关注剖面的拟合精度，更通过 \texttt{compare\_profiles.py} 等工具评估不同拟合方法对下游 LHW 电流驱动剖面的实际影响，将「拟合精度的改善」与「物理计算结果的改善」直接关联。

\section{总结：本文项目的定位}

综合上述三条脉络，本文项目的创新点可归纳为：

\begin{enumerate}[label=\arabic*.]
    \item \textbf{方法论交叉}：将 DeepONet 的「分支--主干」算子学习范式与物理约束损失（PI 思想的推广）相结合，用于托卡马克 $T_e$/$n_e$/$T_i$ 剖面的数据驱动重构，替代 mtanh+spline 参数化拟合链。
    \item \textbf{系统化架构对比}：在同一 DeepONet 框架下实现 SetEncoder（PointNet 式）、BiLSTM、TransformerEncoder 三种编码器，并以纯 ResNet CNN 作为非 DeepONet 对照组，定量回答「分支--主干点积」结构是否为性能关键。
    \item \textbf{物理约束的显式编码}：以单调性、磁轴对称性、光滑性三条剖面先验作为软惩罚，验证物理先验对数据驱动拟合的增益（消融实验）。
    \item \textbf{下游闭环验证}：将拟合剖面接入 ONETWO 输运流程与 LHCD 计算，评估不同拟合方法对电流驱动剖面的实际影响，突破「只比拟合精度、不比物理后果」的局限。
\end{enumerate}

现有文献分别覆盖了神经算子理论（DeepONet/FNO）、物理信息约束（PINN/PI-DeepONet）、聚变剖面重建（RTCAKENN/NN-EFIT）等单个维度，但将三者系统结合、并以四种架构对照 + 物理约束消融 + 下游 LHCD 闭环验证的方式作用于 EAST 实测数据的完整研究尚属空白。本文正是对这一空白的填补。

\section*{参考文献（高相关性文献列表）}

以下为与本文项目相关性最高的文献，按主题分组：

\noindent\textbf{A. 神经算子与物理信息学习（方法核心）}
\begin{enumerate}[label={[\arabic*]}]
    \item Lu L, Jin P, Karniadakis G E. DeepONet: Learning nonlinear operators for identifying differential equations based on the universal approximation theorem of operators. Nature Machine Intelligence, 2021.
    \item Wang S, Wang H, Perdikaris P. Learning the solution operator of parametric PDEs with physics-informed DeepONets. arXiv:2103.10974, 2021.
    \item Karniadakis G E, et al. Physics-informed machine learning. Nature Reviews Physics, 2021.
    \item Gopakumar V, et al. Plasma surrogate modelling using Fourier Neural Operators. arXiv:2311.05967, 2024.
\end{enumerate}

\noindent\textbf{B. 网络架构基础（组件）}
\begin{enumerate}[label={[\arabic*]}]
    \item Qi C R, et al. PointNet: Deep learning on point sets for 3D classification and segmentation. CVPR, 2017.
    \item Vaswani A, et al. Attention is all you need. NeurIPS, 2017.
    \item Hochreiter S, Schmidhuber J. Long short-term memory. Neural Computation, 1997.
\end{enumerate}

\noindent\textbf{C. 剖面拟合与聚变 ML 应用（领域应用）}
\begin{enumerate}[label={[\arabic*]}]
    \item Groebner R J, et al. Progress in quantifying the edge physics of the H-mode regime in DIII-D. Nucl. Fusion 41, 1789 (2001).
    \item Shousha R, et al. Machine learning-based real-time kinetic profile reconstruction in DIII-D. Nucl. Fusion 64, 026006 (2024).
    \item Sun X, et al. Impact of various DIII-D diagnostics on the accuracy of neural network surrogates for kinetic EFIT reconstructions. Nucl. Fusion 64, 086065 (2024).
    \item Lu J, Hu Y, Xiang N, Sun Y. Fast equilibrium reconstruction by deep learning on EAST tokamak. arXiv:2305.12098, 2023.
\end{enumerate}

\noindent\textbf{D. EAST 装置与输运建模（背景）}
\begin{enumerate}[label={[\arabic*]}]
    \item Zang Q, et al. Upgraded multipulse laser and multipoint Thomson scattering diagnostics on EAST. Rev. Sci. Instrum. 82, 063502 (2011).
    \item Meneghini O, et al. Integrated modeling applications for tokamak experiments with OMFIT. Nucl. Fusion 55, 083008 (2015).
\end{enumerate}

\end{document}
